% Trinity S³AI NeurIPS 2026 Paper Template
% Based on NeurIPS 2025 format

\documentclass{article}

% ===== PACKAGES =====
\usepackage[preprint]{neurips_2024}
\usepackage[utf8]{inputenc}
\usepackage[T1]{fontenc}
\usepackage{hyperref}
\usepackage{url}
\usepackage{booktabs}
\usepackage{amsfonts}
\usepackage{nicefrac}
\usepackage{microtype}
\usepackage{xcolor}
\usepackage{amsmath}
\usepackage{amssymb}
\usepackage{amsthm}
\usepackage{algorithm}
\usepackage{algpseudocode}
\usepackage{graphicx}

% ===== METADATA =====
\title{Trinity S³AI: Ternary Sparse AI for Edge Deployment}

\author{
  Dmitrii Vasilev \\
  Trinity Research Collective
}

\begin{document}

\maketitle

\begin{abstract}
We introduce Trinity S³AI (Sparse, Sacred, Scalable Artificial Intelligence), a framework for efficient ternary neural networks. Key innovations: (1) balanced ternary computing \{-1, 0, +1\} achieving 20× memory compression, (2) sacred scaling based on $\phi^2 + \phi^{-2} = 3$, (3) sparse VSA with 90\% sparsity. HSLM-1.95M achieves PPL=125.3 with 533× energy efficiency.
\end{abstract}

\section{Introduction}

\subsection{Trinity Identity}

The golden ratio $\phi = (1+\sqrt{5})/2 \approx 1.618$ satisfies:

\begin{equation}
\phi^2 + \phi^{-2} = 3
\end{equation}

This identity provides the theoretical foundation for sacred scaling.

\section{Method}

\subsection{Ternary Quantization}

\begin{equation}
W_Q[i,j] = \begin{cases}
+1 & \text{if } W[i,j] > \phi^{-1} \sigma \\
0 & \text{if } |W[i,j]| \leq \phi^{-1} \sigma \\
-1 & \text{if } W[i,j] < -\phi^{-1} \sigma
\end{cases}
\end{equation}

\subsection{Sacred Scaling}

\begin{equation}
\sigma_{\text{sacred}} = d^{-\phi^{-3}} = d^{-0.236}
\end{equation}

\section{Results}

\begin{table}[h]
\centering
\caption{Perplexity on TinyStories}
\begin{tabular}{lcc}
\toprule
Method & PPL & CI95 \\
\midrule
Standard & 128.7 & [125.9, 131.5] \\
\textbf{Sacred} & \textbf{125.3} & \textbf{[123.1, 127.5]} \\
\bottomrule
\end{tabular}
\end{table}

\section{Conclusion}

Trinity S³AI achieves 125.3 PPL with 20× compression and 533× energy efficiency.

\end{document}
